\documentclass[12pt, a4paper]{article}

\usepackage[utf8]{inputenc}
\usepackage[T1]{fontenc}
\usepackage[french]{babel}
\usepackage{amsmath, amssymb, amsthm}
\usepackage{geometry}
\usepackage{hyperref}
\usepackage{booktabs}
\usepackage{array}
\usepackage{tikz}
\usetikzlibrary{arrows.meta, positioning, calc, shapes.geometric}
\usepackage{xcolor}
\usepackage{float}
\usepackage{multirow}

\geometry{margin=2.5cm}

\definecolor{bleu}{RGB}{30, 100, 200}
\definecolor{rouge}{RGB}{200, 50, 50}
\definecolor{vert}{RGB}{30, 150, 80}
\definecolor{orange}{RGB}{210, 120, 0}
\definecolor{violet}{RGB}{130, 50, 180}
\definecolor{gris}{RGB}{100, 100, 100}

\hypersetup{colorlinks=true, linkcolor=bleu, urlcolor=bleu}

\newtheorem{definition}{Définition}
\newtheorem{propriete}{Propriété}
\newtheorem{theoreme}{Théorème}
\newtheorem{lemme}{Lemme}
\newcommand{\dist}[2]{d\!\left(#1,\,#2\right)}

\title{
    \textbf{Arbre de Steiner euclidien pour 5 points}\\[0.4em]
    \large Fonctionnement complet de l'algorithme —\\
    du clic utilisateur au résultat affiché
}
\author{Projet Algorithmes 2}
\date{\today}

% ------------------------------------------------------------------
\begin{document}
\maketitle
\tableofcontents
\newpage

% ==================================================================
\section{Rappels fondamentaux}
% ==================================================================

\subsection{Propriété des 120°}

Dans tout arbre de Steiner euclidien optimal, les arêtes forment des angles
de \textbf{exactement 120°} en chaque point de Steiner.

\subsection{Nombre maximal de points de Steiner pour $n = 5$}

\begin{theoreme}
Pour 5 terminaux, un arbre de Steiner optimal contient \textbf{au plus 3 points de Steiner}.
\end{theoreme}

\begin{proof}
Un arbre sur $n + k$ nœuds possède $n + k - 1$ arêtes.
Chaque point de Steiner est de degré exactement 3 dans la solution optimale.
La somme des degrés vaut $2(n + k - 1)$.
Les $n$ terminaux ont degré $\geq 1$, les $k$ Steiner contribuent $3k$ :
\[
    2(n + k - 1) \;\geq\; n + 3k
    \;\Longrightarrow\; k \;\leq\; n - 2.
\]
Pour $n = 5$ : $k \leq 3$.
\end{proof}

\subsection{Topologies possibles}

Un arbre de Steiner sur 5 terminaux peut donc avoir 0, 1, 2 ou 3 points de Steiner.
Le nombre de topologies à énumérer est résumé dans le tableau suivant :

\begin{center}
\begin{tabular}{clcc}
\toprule
$k$ & \textbf{Structure} & \textbf{Nb. topologies} & \textbf{Méthode} \\
\midrule
0 & MST des 5 terminaux & 1 & Prim \\
1 & Fermat(triplet) + 2 terminaux raccrochés & $\sim$150 & Weiszfeld \\
2 & Type I : $S_0$–$S_1$ directs + queue & $\sim$60 & Weiszfeld alterné \\
  & Type II : $T_\text{mid}$ entre $S_0$ et $S_1$ & 15 & Fermat exact \\
3 & Chaîne $S_0$–$S_1$–$S_2$ & 15 & Weiszfeld alterné \\
\midrule
& \textbf{Total} & $\mathbf{\sim 241}$ & \\
\bottomrule
\end{tabular}
\end{center}

% ==================================================================
\section{Famille 0 — Arbre couvrant minimal (MST)}
% ==================================================================

Le MST constitue la référence initiale : toute topologie avec des points de Steiner
doit produire un arbre \textbf{strictement plus court} pour être retenue.

L'algorithme de Prim s'applique sur les 5 terminaux :

\begin{align}
    &\text{inMST} = \{T_1\}, \quad \text{minDist}[i] = +\infty, \quad \text{parent}[i] = -1 \nonumber\\
    &\text{Répéter 4 fois :} \nonumber\\
    &\quad u \leftarrow \arg\min_{i \notin \text{inMST}} \text{minDist}[i] \nonumber\\
    &\quad \text{inMST} \leftarrow \text{inMST} \cup \{u\} \nonumber\\
    &\quad \forall v \notin \text{inMST} : \quad
       \text{si } \dist{T_u}{T_v} < \text{minDist}[v] \;\text{ alors mettre à jour} \nonumber
\end{align}

\textbf{Propriété :} $\ell_\text{MST} \leq 2 \times \ell_\text{Steiner}^*$
(le MST est une 2-approximation de l'arbre de Steiner optimal).

% ==================================================================
\section{Famille 1 — Un point de Steiner ($k = 1$)}
% ==================================================================

\subsection{Construction des topologies}

On choisit un triplet de 3 terminaux parmi 5 : $\binom{5}{3} = 10$ triplets.
Pour chaque triplet $\{T_{i_0}, T_{i_1}, T_{i_2}\}$, les deux terminaux restants
$\ell_0$ et $\ell_1$ doivent être raccordés à l'arbre partiel
$\{S_0, T_{i_0}, T_{i_1}, T_{i_2}\}$ sans créer de cycle.

\medskip

$S_0$ est déjà de degré 3 (plein), donc $\ell_0$ et $\ell_1$ s'accrochent aux
\emph{terminaux} du triplet ou l'un à l'autre (chaîne). Cela donne :

\begin{itemize}
    \item $\ell_0 \to$ nœud du triplet (3 choix), $\ell_1 \to$ nœud du triplet ou $\ell_0$ (4 choix) : $3 \times 4 = 12$ topologies par triplet.
    \item $\ell_1 \to$ nœud du triplet (3 choix), $\ell_0 \to \ell_1$ (chaîne inversée) : 3 topologies par triplet.
\end{itemize}

\textbf{Total :} $10 \times (12 + 3) = 150$ topologies pour $k = 1$.

\subsection{Calcul du point de Steiner}

Le point $S_0$ est le \textbf{point de Fermat-Torricelli} du triplet, calculé par
itérations de Weiszfeld (voir section~4 du document 4-points) :

\[
    S_0^{(k+1)} \;=\;
    \frac{\dfrac{T_{i_0}}{\dist{S_0^{(k)}}{T_{i_0}}}
        + \dfrac{T_{i_1}}{\dist{S_0^{(k)}}{T_{i_1}}}
        + \dfrac{T_{i_2}}{\dist{S_0^{(k)}}{T_{i_2}}}}
         {\dfrac{1}{\dist{S_0^{(k)}}{T_{i_0}}}
        + \dfrac{1}{\dist{S_0^{(k)}}{T_{i_1}}}
        + \dfrac{1}{\dist{S_0^{(k)}}{T_{i_2}}}}
\]

Longueur totale de la topologie :
\[
    \ell \;=\;
    \dist{S_0}{T_{i_0}} + \dist{S_0}{T_{i_1}} + \dist{S_0}{T_{i_2}}
    + \dist{\ell_0}{\text{attach}_0} + \dist{\ell_1}{\text{attach}_1}.
\]

% ==================================================================
\section{Famille 2 — Deux points de Steiner ($k = 2$)}
% ==================================================================

Deux structures topologiques distinctes existent pour $k = 2$ et $n = 5$.

% ------------------------------------------------------------------
\subsection{Type I — $S_0$ et $S_1$ directement connectés}
% ------------------------------------------------------------------

\subsubsection{Structure}

Un terminal «~queue~» ($T_\text{tail}$) se raccorde à l'un des 4 terminaux du cœur.
Les 4 terminaux du cœur sont répartis en deux paires $(T_a, T_b)$ et $(T_c, T_d)$ :

\[
    S_0 \;\leftrightarrow\; T_a,\; T_b,\; S_1
    \qquad
    S_1 \;\leftrightarrow\; T_c,\; T_d,\; S_0
    \qquad
    T_\text{tail} \;\to\; T_\text{attach} \in \{T_a, T_b, T_c, T_d\}
\]

\begin{center}
\begin{tikzpicture}[
    pt/.style={circle, fill=bleu!80, minimum size=6pt, inner sep=0pt},
    st/.style={circle, fill=rouge!80, minimum size=6pt, inner sep=0pt},
    edge/.style={thick, gris}
]
    \node[pt, label=above:$T_a$] (Ta) at (0,1.5) {};
    \node[pt, label=below:$T_b$] (Tb) at (0,0) {};
    \node[pt, label=above:$T_c$] (Tc) at (4,1.5) {};
    \node[pt, label=below:$T_d$] (Td) at (4,0) {};
    \node[pt, label=below:$T_\text{tail}$] (Tt) at (2,-1.2) {};
    \node[st, label=above:$S_0$] (S0) at (1.3,0.75) {};
    \node[st, label=above:$S_1$] (S1) at (2.7,0.75) {};

    \draw[edge] (S0) -- (Ta);
    \draw[edge] (S0) -- (Tb);
    \draw[edge] (S0) -- (S1);
    \draw[edge] (S1) -- (Tc);
    \draw[edge] (S1) -- (Td);
    \draw[edge, dashed] (Tt) -- (Td) node[midway, right, font=\footnotesize]{attach};
\end{tikzpicture}
\end{center}

\subsubsection{Comptage}

5 choix pour $T_\text{tail}$ $\times$ 3 partitions des 4 terminaux restants $\times$
4 choix d'attache $= 60$ topologies.

\subsubsection{Optimisation}

Les positions de $S_0$ et $S_1$ sont optimisées par \textbf{Weiszfeld alterné}
(même principe que pour $k = 2$ dans le cas 4 points) :

\[
    S_0^{(k+1)} = \text{Weber}\!\left(T_a,\, T_b,\, S_1^{(k)}\right)
    \qquad
    S_1^{(k+1)} = \text{Weber}\!\left(T_c,\, T_d,\, S_0^{(k+1)}\right)
\]

où $\text{Weber}(A, B, C)$ désigne la mise à jour de Weiszfeld :
\[
    \text{Weber}(A, B, C) \Big|_{F^{(k)}}
    \;=\;
    \frac{A/\dist{F^{(k)}}{A} \;+\; B/\dist{F^{(k)}}{B} \;+\; C/\dist{F^{(k)}}{C}}
         {1/\dist{F^{(k)}}{A} \;+\; 1/\dist{F^{(k)}}{B} \;+\; 1/\dist{F^{(k)}}{C}}.
\]

Longueur totale :
\[
    \ell \;=\;
    \dist{S_0}{T_a} + \dist{S_0}{T_b} + \dist{S_0}{S_1}
    + \dist{S_1}{T_c} + \dist{S_1}{T_d}
    + \dist{T_\text{tail}}{T_\text{attach}}.
\]

% ------------------------------------------------------------------
\subsection{Type II — Terminal $T_\text{mid}$ entre $S_0$ et $S_1$}
% ------------------------------------------------------------------

\subsubsection{Structure}

Un terminal $T_\text{mid}$ est le nœud central : il est connecté à la fois à
$S_0$ et à $S_1$. Les 4 autres terminaux sont répartis en deux paires :

\[
    S_0 \;\leftrightarrow\; T_a,\; T_b,\; T_\text{mid}
    \qquad
    S_1 \;\leftrightarrow\; T_c,\; T_d,\; T_\text{mid}
\]

\begin{center}
\begin{tikzpicture}[
    pt/.style={circle, fill=bleu!80, minimum size=6pt, inner sep=0pt},
    st/.style={circle, fill=rouge!80, minimum size=6pt, inner sep=0pt},
    edge/.style={thick, gris}
]
    \node[pt, label=above:$T_a$] (Ta) at (0,1.5) {};
    \node[pt, label=below:$T_b$] (Tb) at (0,0) {};
    \node[pt, label=above:$T_\text{mid}$] (Tm) at (2,0.75) {};
    \node[pt, label=above:$T_c$] (Tc) at (4,1.5) {};
    \node[pt, label=below:$T_d$] (Td) at (4,0) {};
    \node[st, label=below:$S_0$] (S0) at (0.9,0.75) {};
    \node[st, label=below:$S_1$] (S1) at (3.1,0.75) {};

    \draw[edge] (S0) -- (Ta);
    \draw[edge] (S0) -- (Tb);
    \draw[edge] (S0) -- (Tm);
    \draw[edge] (S1) -- (Tm);
    \draw[edge] (S1) -- (Tc);
    \draw[edge] (S1) -- (Td);
\end{tikzpicture}
\end{center}

\subsubsection{Propriété clé}

$S_0$ et $S_1$ sont \textbf{indépendants} : ils ne partagent aucune arête entre eux,
uniquement le terminal fixe $T_\text{mid}$. Donc chacun est le point de Fermat exact
de son triplet de terminaux :
\[
    S_0 = \text{Fermat}(T_a,\, T_b,\, T_\text{mid})
    \qquad
    S_1 = \text{Fermat}(T_c,\, T_d,\, T_\text{mid}).
\]

\subsubsection{Comptage}

5 choix pour $T_\text{mid}$ $\times$ 3 partitions des 4 autres terminaux $= 15$ topologies.

% ==================================================================
\section{Famille 3 — Trois points de Steiner ($k = 3$)}
% ==================================================================

\subsection{Unicité de la structure}

\begin{theoreme}
Pour $n = 5$ terminaux et $k = 3$ points de Steiner, la seule topologie arborescente
valide est la \textbf{chaîne} $S_0 - S_1 - S_2$.
\end{theoreme}

\begin{proof}
L'arbre a $5 + 3 = 8$ nœuds et 7 arêtes.
Chaque $S_i$ est de degré 3. La somme des degrés vaut $3 \times 3 + 5 \times d_j = 14$,
donc les 5 terminaux ont degré total $14 - 9 = 5$, soit degré moyen 1.
L'unique structure arborescente où 3 nœuds internes ont degré 3 est la chaîne
$S_0 - S_1 - S_2$ où $S_0$ et $S_2$ portent chacun 2 terminaux feuilles
et $S_1$ porte 1 terminal feuille.
\end{proof}

\subsection{Structure de la chaîne}

\[
    S_0 \;\leftrightarrow\; T_a,\; T_b,\; S_1
    \qquad
    S_1 \;\leftrightarrow\; S_0,\; T_c,\; S_2
    \qquad
    S_2 \;\leftrightarrow\; S_1,\; T_d,\; T_e
\]

\begin{center}
\begin{tikzpicture}[
    pt/.style={circle, fill=bleu!80, minimum size=6pt, inner sep=0pt},
    st/.style={circle, fill=rouge!80, minimum size=6pt, inner sep=0pt},
    edge/.style={thick, gris}
]
    \node[pt, label=above left:$T_a$] (Ta) at (-1, 1.2) {};
    \node[pt, label=below left:$T_b$] (Tb) at (-1,-0.2) {};
    \node[st, label=above:$S_0$] (S0) at (0.2, 0.5) {};
    \node[pt, label=above:$T_c$] (Tc) at (2, 1.2) {};
    \node[st, label=above:$S_1$] (S1) at (2, 0.5) {};
    \node[st, label=above:$S_2$] (S2) at (3.8, 0.5) {};
    \node[pt, label=above right:$T_d$] (Td) at (5, 1.2) {};
    \node[pt, label=below right:$T_e$] (Te) at (5,-0.2) {};

    \draw[edge] (S0)--(Ta); \draw[edge] (S0)--(Tb);
    \draw[edge] (S0)--(S1);
    \draw[edge] (S1)--(Tc); \draw[edge] (S1)--(S2);
    \draw[edge] (S2)--(Td); \draw[edge] (S2)--(Te);
\end{tikzpicture}
\end{center}

\subsection{Comptage des topologies}

\begin{itemize}
    \item 5 choix pour $T_c$ (le terminal «~milieu~» connecté à $S_1$).
    \item $\binom{4}{2}/2 = 3$ façons de répartir les 4 terminaux restants en deux paires
          non ordonnées $(T_a, T_b)$ et $(T_d, T_e)$ (division par 2 pour la symétrie $S_0 \leftrightarrow S_2$).
\end{itemize}

\textbf{Total :} $5 \times 3 = 15$ topologies.

\subsection{Optimisation des trois Steiner — Weiszfeld à 3 blocs}

On minimise simultanément :
\[
    f(S_0, S_1, S_2) =
    \dist{S_0}{T_a} + \dist{S_0}{T_b} + \dist{S_0}{S_1}
    + \dist{S_1}{T_c} + \dist{S_1}{S_2}
    + \dist{S_2}{T_d} + \dist{S_2}{T_e}.
\]

Les trois mises à jour de Weiszfeld sont effectuées en cascade à chaque itération :

\begin{align}
    S_0^{(k+1)} &= \text{Weber}\!\left(T_a,\; T_b,\; S_1^{(k)}\right)
    \Big|_{S_0^{(k)}}
    \label{eq:s0} \\[8pt]
    S_1^{(k+1)} &= \text{Weber}\!\left(S_0^{(k+1)},\; T_c,\; S_2^{(k)}\right)
    \Big|_{S_1^{(k)}}
    \label{eq:s1} \\[8pt]
    S_2^{(k+1)} &= \text{Weber}\!\left(S_1^{(k+1)},\; T_d,\; T_e\right)
    \Big|_{S_2^{(k)}}
    \label{eq:s2}
\end{align}

où la notation $\text{Weber}(A, B, C)\big|_{F^{(k)}}$ désigne :
\[
    \text{Weber}(A, B, C)\Big|_{F^{(k)}}
    \;=\;
    \frac{\dfrac{A}{\dist{F^{(k)}}{A}}
        + \dfrac{B}{\dist{F^{(k)}}{B}}
        + \dfrac{C}{\dist{F^{(k)}}{C}}}
         {\dfrac{1}{\dist{F^{(k)}}{A}}
        + \dfrac{1}{\dist{F^{(k)}}{B}}
        + \dfrac{1}{\dist{F^{(k)}}{C}}}
\]

\textbf{Critère d'arrêt :}
\[
    \max\!\left(
        \|S_0^{(k+1)} - S_0^{(k)}\|,\;
        \|S_1^{(k+1)} - S_1^{(k)}\|,\;
        \|S_2^{(k+1)} - S_2^{(k)}\|
    \right) \;<\; 10^{-10}
\]

\subsection{Initialisations multiples}

Pour la chaîne $S_0$–$S_1$–$S_2$, trois positions de départ sont testées :

\begin{center}
\begin{tabular}{ll}
\toprule
\textbf{Init} & \textbf{Positions initiales $(S_0^{(0)},\; S_1^{(0)},\; S_2^{(0)})$} \\
\midrule
1 & Milieux $\overline{T_a T_b}$, centroïde global, milieu $\overline{T_d T_e}$ \\
2 & Fermat$(T_a, T_b, T_c)$, centroïde global, Fermat$(T_d, T_e, T_c)$ \\
3 & $\tfrac{1}{2}(T_a + G)$,\; $\tfrac{1}{2}(T_c + G)$,\; $\tfrac{1}{2}(T_d + G)$ où $G$ = centroïde \\
\bottomrule
\end{tabular}
\end{center}

% ==================================================================
\section{Validation géométrique}
% ==================================================================

Après chaque optimisation ($k = 2$ ou $k = 3$), la solution est rejetée si l'une
des conditions suivantes est vérifiée :

\begin{enumerate}
    \item \textbf{Steiner confondu avec un terminal :}
          $\exists\, i,\, j : \dist{S_i}{T_j} < \varepsilon$
    \item \textbf{Deux Steiner confondus :}
          $\exists\, i \neq j : \dist{S_i}{S_j} < \varepsilon$
\end{enumerate}

Le seuil adaptatif est $\varepsilon = \max(10^{-6},\; 10^{-3} \times \Delta)$
où $\Delta = \max_{i < j}\dist{T_i}{T_j}$.

% ==================================================================
\section{Sélection finale}
% ==================================================================

L'algorithme maintient un unique «~meilleur~» résultat initialisé avec le MST.
Après chaque topologie valide évaluée, si sa longueur est inférieure :

\[
    \ell(\text{candidat}) < \ell_\text{best}
    \;\Longrightarrow\;
    \text{best} \leftarrow \text{candidat}, \quad
    \ell_\text{best} \leftarrow \ell(\text{candidat}).
\]

Le résultat final est \textbf{toujours valide} car le MST constitue une borne supérieure.

% ==================================================================
\section{Flux de données complet}
% ==================================================================

\begin{center}
\begin{tikzpicture}[
    step/.style={draw, rounded corners, fill=#1!10, draw=#1!60,
                 minimum width=7.5cm, minimum height=0.75cm, font=\small},
    arr/.style={-{Stealth}, thick}
]
    \node[step=bleu]   (A)  {Utilisateur place 5 points sur le canvas};
    \node[step=bleu,  below=0.3cm of A]  (B)  {AppComponent.onPointAdded() $\times$ 5};
    \node[step=bleu,  below=0.3cm of B]  (C)  {SteinerService.solve() — POST JSON};
    \node[step=rouge, below=0.3cm of C]  (D)  {SteinerController.solve()};
    \node[step=rouge, below=0.3cm of D]  (E)  {SteinerTreeService.solveForFivePoints()};
    \node[step=vert,  below=0.3cm of E]  (F0) {Évaluation MST — initialise best};
    \node[step=vert,  below=0.3cm of F0] (F1) {$k=1$ : 150 topologies — Fermat-Weiszfeld};
    \node[step=orange,below=0.3cm of F1] (F2a){$k=2$ Type I : 60 topologies — Weiszfeld alterné};
    \node[step=orange,below=0.3cm of F2a](F2b){$k=2$ Type II : 15 topologies — Fermat exact};
    \node[step=violet,below=0.3cm of F2b](F3) {$k=3$ : 15 topologies — Weiszfeld 3 blocs};
    \node[step=rouge, below=0.3cm of F3] (G)  {Sélection du minimum — SteinerResult};
    \node[step=bleu,  below=0.3cm of G]  (H)  {Affichage sur le canvas Angular};

    \foreach \x/\y in {A/B, B/C, C/D, D/E, E/F0, F0/F1, F1/F2a, F2a/F2b, F2b/F3, F3/G, G/H}
        \draw[arr] (\x) -- (\y);
\end{tikzpicture}
\end{center}

% ==================================================================
\section{Exemple numérique — pentagone régulier}
% ==================================================================

Soit les 5 sommets d'un pentagone régulier de rayon $R = 1$ :
\[
    T_i = \left(\cos\frac{2\pi i}{5},\; \sin\frac{2\pi i}{5}\right), \quad i = 0, \ldots, 4.
\]

\textbf{MST :} longueur $\approx 5 \times 2\sin(\pi/5) / \text{quelques arêtes}
\approx 4{,}702$.

\textbf{Résultat optimal :} la topologie $k = 3$ en chaîne donne une longueur
$\approx 4{,}509$, meilleure que le MST et toutes les topologies $k \leq 2$.

% ==================================================================
\section{Complexité}
% ==================================================================

\begin{center}
\begin{tabular}{lccc}
\toprule
\textbf{Phase} & \textbf{Topologies} & \textbf{Coût par topologie} & \textbf{Total} \\
\midrule
MST (Prim)           & 1   & $O(n^2)$              & $O(25)$ \\
$k=1$ (Weiszfeld)    & 150 & $O(\text{iter})$      & $O(150 \cdot \text{iter})$ \\
$k=2$ Type I         & 60  & $O(4 \cdot \text{iter})$ & $O(240 \cdot \text{iter})$ \\
$k=2$ Type II        & 15  & $O(\text{iter})$      & $O(15 \cdot \text{iter})$ \\
$k=3$                & 15  & $O(3 \cdot \text{iter})$ & $O(45 \cdot \text{iter})$ \\
\midrule
\textbf{Total}       & $\sim\!241$ & — & $O(450 \cdot \text{iter})$ \\
\bottomrule
\end{tabular}
\end{center}

\medskip

Avec iter $\leq 50\,000$ et une convergence typique en quelques centaines d'itérations,
l'algorithme reste très rapide en pratique pour $n = 5$.

\end{document}
